\scanpage{91}
\uline{NB} De façon générale, s'intéressant aux gpes $H$ (finis) dans lesquels chaque él. est d'ordre \deleted{fini} divisant un entier $N$ fixé ($N=6$ pour le gpe tétraédral, $N=12$ pour celui ${\unsure{G|\mathfrak S|\mathcal S}}_4$ de l'octaèdre, $N=30$ pour celui $A_5$ de l'icosaèdre) on est amené à remplacer $\Pi$ par le quotient correspondant $\Pi(N)$, $\Omega$ par $\Omega(\deleted{\unsure{6}}N)=\mathrm{Aut}(\Pi(N))$, \deleted{et} $G_3\underset{1/2}{\cdot}|\unsure{E|\mathcal E|\mathbb E}|$ par $G_3\underset{1/2}{\cdot}|E|(N)$, quotient de \deleted{$G_3\underset{1/2}{\cdot}(\gamma(N)\times\gamma(N))$}
\[
\left\{\begin{aligned}
&G_3\underset{1/2}{\cdot}(\gamma(N)\times\unsure{\gamma|\sigma}(N))^{(*\deleted{\unsure{\gamma}}G_3/\pi_{01})}\\
&\text{où}\qquad\gamma(N)=(\mathbb Z/N\mathbb Z)^*
\end{aligned}\right.
\]
\[
\left[\,\gamma(N)\simeq
\left|\begin{aligned}
&\pm1 &&\text{pour }A_{\deleted{\unsure{5}}4}\\
&(\pm1)\times(\pm1)&&\text{pour }G_{\deleted{\unsure{5}}4}\\
&(\pm1)\times(\mathbb Z/4)&&\text{pour }A_5\\
&(\pm1)\times(\pm1)\times(\mathbb Z/4)&&\text{pour }G_5\qquad\text{etc.}
\end{aligned}\right.\right.
\]
% archive sheet number: 90
